% Elaborado por Fabian Gomez
% Version 1.0

\section{Appendices}\label{sec:appendices}

Esta sección agrupa anexos de apoyo para el SyRS/SRS final del Rover Olympus. Incluye (i) un glosario consolidado (sinónimos/abreviaturas y términos de dominio) y (ii) un resumen de historial de revisiones, depurando inconsistencias entre los borradores main/main1/main3 y dejando un criterio claro de baseline para el documento final. \cite{omg_sysml_mbse, nasa_cm, segoldmine_ppi}  
Los anexos se mantienen concisos: el propósito es facilitar lectura y auditoría académica, sin duplicar contenido ``maestro'' que ya está en secciones 1--13 (requisitos, V\&V, trazabilidad, arquitectura). \cite{omg_sysml_mbse, nasa_cm, segoldmine_ppi}  

\subsection{Traceability Matrix (SyRS $\rightarrow$ Subsystem $\rightarrow$ V\&V)}
La siguiente tabla resume la trazabilidad de los requerimientos de sistema (SyRS) hacia sus derivados por subsistema y el método de verificación primaria (Test, Analysis, Demonstration, Inspection).

\begin{table}[H]
\centering
\small
\begin{tabularx}{\textwidth}{@{} l p{5.5cm} l l @{}}
\toprule
\textbf{ID SyRS} & \textbf{Descripción Resumida} & \textbf{Derivado (Hijo)} & \textbf{Método V\&V} \\ \midrule
RF-001 & Percepción Visual & GNC-REQ-001 & Demonstration \\
RF-002 & Clasificación de Terreno & GNC-REQ-002 & Analysis \\
RF-003 & Evitación de Obstáculos & TRAC-REQ-001, GNC-REQ-002 & Test \\
RF-004 & Compensación de Slip & GNC-REQ-003 & Test \\
RF-005 & Fault Stop Seguro & PROP-REQ-002 & Demonstration \\
RF-006 & Integridad de Datos & COMM-REQ-001, 004 & Analysis \\
RNF-001 & Latencia de Procesamiento $\le$ 2s & CDH-REQ-001 & Test \\
RNF-002 & Resiliencia de Recuperación & ConOps/SW-logic & Demonstration \\
USA-REQ-002 & Despliegue Operativo $\le$ 15 min & PKG-REQ-001 & Demonstration \\
SEC-REQ-001 & Autenticación de Comandos & COMM-REQ-001 & Analysis \\
INF-REQ-001 & Persistencia de Logs & CDH-REQ-002 & Inspection \\
PHY-REQ-001 & Dimensiones 480x340x320 mm & STR-REQ-001 & Inspection \\
\bottomrule
\end{tabularx}
\caption{Matriz de Trazabilidad Maestra del Rover Olympus.}
\end{table}

\subsection{TBD Register (Plan de Resolución)}
Conforme a ISO/IEC/IEEE 29148, este registro identifica los parámetros técnicos aún no cuantificados, asigna un responsable y establece un criterio para su resolución antes de la revisión final de V\&V.

\begin{table}[H]
\centering
\scriptsize
\begin{tabularx}{\textwidth}{| l | X | l | l | l |}
\hline
\textbf{TBD-ID} & \textbf{Descripción del Parámetro} & \textbf{Requisito/Elemento} & \textbf{Responsable} & \textbf{Criterio de Resolución} \\ \hline
TBD-OP-01 & Umbrales de severidad para contingencias (\texttt{severity\_allows\_retreat}) & §7.3.8 (SM), s09a & Líder GNC & Análisis de simulación en entorno volcánico. \\ \hline
TBD-OP-02 & Política de reintentos de enlace (\texttt{retry\_allowed}) & §7.3.8 (SM), s09a & Líder COMM & Caracterización de latencia UHF en campo. \\ \hline
TBD-ENV-01 & Nivel de protección antipartículas (IP4X o superior) & STR-REQ-002, SyRS-100 & Líder STR & Pruebas de infiltración con regolito simulado. \\ \hline
TBD-PER-01 & Baseline de referencia para mejora EKF (sin EKF) & SyRS-013, §11.1 & Líder GNC & Definición de ground-truth en sitio análogo. \\ \hline
TBD-DATA-01 & Umbral de almacenamiento crítico (\texttt{storage\_near\_full}) & CDH-REQ-002, §11.1 & Líder C\&DH & Capacidad efectiva de microSD vs bitrate TM. \\ \hline
TBD-LOG-01 & Formato y versión de datasets de misión & §9.5.14 (INF-REQ) & Líder C\&DH & Estandarización conforme a ELANav API. \\ \hline
\end{tabularx}
\caption{Registro maestro de parámetros pendientes (TBD) y plan de resolución.}
\end{table}

\subsection{Glossary}

Este glosario consolida términos y abreviaturas utilizados en el SyRS/SRS final. Para evitar redundancias, se considera complementario a la Sección 3 (Terms, Definitions, Acronyms), que contiene las definiciones operativas principales. \cite{omg_sysml_mbse, nasa_cm, segoldmine_ppi}  

\textbf{Términos clave (consolidación):}
\begin{itemize}
    \item \textbf{ELANav:} Embedded Library for Autonomous Navigation (biblioteca embebida para navegación autónoma). \cite{omg_sysml_mbse, nasa_cm}  
    \item \textbf{Rover Olympus:} Rover académico tipo exploración, plataforma de validación de ELANav (7 subsistemas). \cite{omg_sysml_mbse, nasa_cm, segoldmine_ppi}  
    \item \textbf{UC‑01:} Exploración de superficie (caso núcleo); navegación autónoma continua con percepción y evitación. \cite{segoldmine_ppi, nasa_cm}  
    \item \textbf{Sitio análogo:} Terreno volcánico / regolito simulado usado para validación (pendientes hasta 20°, obstáculos, transiciones de suelo). \cite{nasa_cm, segoldmine_ppi}  
    \item \textbf{Slip ratio / deslizamiento:} Indicador de diferencia entre avance teórico y real; usado como disparador de contingencia con umbral slip $>$ 20\%. \cite{omg_sysml_mbse, nasa_cm, segoldmine_ppi}  
    \item \textbf{GCS / Estación terrena:} Sistema externo operado por Rover Planner para uplink/downlink. \cite{nasa_cm, segoldmine_ppi}  
    \item \textbf{CSP:} CubeSat Space Protocol (encapsulamiento TM/CMD). \cite{omg_sysml_mbse, nasa_cm, segoldmine_ppi}  
    \item \textbf{V\&V:} Verification \& Validation, con métodos T/A/D/I (Test/Analysis/Demonstration/Inspection). \cite{omg_sysml_mbse, nasa_cm, segoldmine_ppi}  
    \item \textbf{BDD/IBD/PBS/N2/ASM:} Artefactos de arquitectura (SysML y descomposición física) usados para trazabilidad y soporte de integración. \cite{omg_sysml_mbse, segoldmine_ppi}  
\end{itemize}

\subsection{Revision History}

Este SyRS/SRS final es el resultado de un proceso iterativo de borradores (main, main1, main3) que fueron generados en distintas iteraciones y niveles de detalle. En lugar de mantener historias de revisión divergentes por archivo, esta sección consolida el historial relevante y establece un criterio único de baseline para el documento final. \cite{omg_sysml_mbse, nasa_cm, segoldmine_ppi}  

\textbf{Historial consolidado (síntesis):}
\begin{itemize}
    \item \textbf{Iteración inicial:} Creación del documento base y establecimiento del foco UC‑01 como caso núcleo de validación. \cite{nasa_cm, segoldmine_ppi}  
    \item \textbf{Reestructuración:} Adopción de estructura alineada a ISO/IEC/IEEE 29148 y guía NASA SE (reorganización de secciones, formalización de V\&V y trazabilidad). \cite{nasa_cm, segoldmine_ppi, omg_sysml_mbse}  
    \item \textbf{Auditoría e integración:} Consolidación de arquitectura lógica/física (7 subsistemas), diferenciación PBS vs BDD, inclusión de PBS/ASM y matriz N2, y fortalecimiento del plan V\&V con criterios cuantitativos. \cite{nasa_cm, segoldmine_ppi}  
    \item \textbf{Versión extendida:} Incorporación de detalle adicional de integración y apéndices ampliados de ConOps/Casos de Uso (main3), usados como insumo para completar secciones de main.pdf. \cite{segoldmine_ppi}  
    \item \textbf{Consolidación final (esta versión):} Unificación, depuración y normalización de contenido entre main/main1/main3, eliminando redundancias y fijando ``definición maestra'' de requisitos en Sección 11; referencias cruzadas por ID en el resto del documento; y ajuste realista del UC‑05 hacia ``Safe Deactivation/Passivation'' documentado como [EXTERNO] en ConOps/Scope. \cite{omg_sysml_mbse, nasa_cm, segoldmine_ppi}  
\end{itemize}

\textbf{Nota (SE) sobre inconsistencias depuradas:}
\begin{itemize}
    \item En main.pdf aparecía texto residual (``TBD'' en algunas secciones y referencias a una ``versión 3.15''); para este SRS final se adopta el control de revisiones más consistente observado en main1/main3 (v3.12 como referencia interna de borrador) y se elimina la numeración residual no sustentada. \cite{omg_sysml_mbse, nasa_cm, segoldmine_ppi}  
\end{itemize}
